\section{Involution of leaves in \(4\) representations}

The involution \(C^*\mapsto-C^*\) organizes four distinct responses:
\begin{equation}
\label{rm:eq:modular-sheet-parities}
\begin{aligned}
S_{\rm bi}(A,-C^*)&=S_{\rm bi}(A,C^*),\\
\Gamma_{6\to8}(A,-C^*)&=-\Gamma_{6\to8}(A,C^*),\\
(\widehat\mu,\widehat\varepsilon)&\longleftrightarrow
(\widehat\varepsilon,\widehat\mu),\\
\widehat c&\longmapsto\widehat c,
\qquad D_A\longmapsto D_A.
\end{aligned}
\end{equation}
The bilayer information is even; the areal orientation is odd; the two
constitutive quantities form a doublet, and the determinants remain fixed. This parity
table precedes any interpretation as a physical symmetry.

\section{Involutive rational transformation in the CKM parameter space}

Let
\begin{equation}
\label{rm:eq:modular-Mckm}
M_{\rm CKM}=
\begin{pmatrix}
2&-5&28\\
0&7&-25/2\\
1&-20&-2/3
\end{pmatrix},
\qquad
\det M_{\rm CKM}=-\frac{3857}{6}.
\end{equation}
The inverse exists on \(\mathbb Q\) and is
\begin{equation}
\label{rm:eq:modular-Mckm-inverse}
M_{\rm CKM}^{-1}=
\begin{pmatrix}
1528/3857&3380/3857&801/3857\\
75/3857&176/3857&-150/3857\\
6/551&-30/551&-12/551
\end{pmatrix}.
\end{equation}
With
\begin{equation}
\label{rm:eq:modular-ckm-coordinates}
h=\frac{C^*}{6},
\qquad
\Delta=\frac{180}{\pi}\Delta_4,
\qquad
\boldsymbol\theta=M_{\rm CKM}(A,h,\Delta)^T,
\end{equation}
the odd coordinate enters in the direction
\begin{equation}
\label{rm:eq:modular-ch}
c_h=\begin{pmatrix}-5\\7\\-20\end{pmatrix}.
\end{equation}

If \(c_A=(2,0,1)^T\) and
\(c_\Delta=(28,-25/2,-2/3)^T\), the three columns
\((c_A,c_h,c_\Delta)\) are precisely the basis transported by
\(M_{\rm CKM}\). The second row of
\cref{rm:eq:modular-Mckm-inverse} is the covector \(\ell_h\) that extracts the
odd coordinate:
\begin{equation}
\label{rm:eq:modular-dual-coordinate-check}
\ell_hc_A=0,
\qquad
\ell_hc_h=1,
\qquad
\ell_hc_\Delta=0.
\end{equation}

\begin{theorem}[Reflexion CKM inducida by the sheet]
\label{rm:thm:modular-ckm-reflection}
Let
\begin{equation}
\label{rm:eq:modular-Dh-lh}
D_h=\operatorname{diag}(1,-1,1),
\qquad
\ell_h=\frac1{3857}\begin{pmatrix}75&176&-150\end{pmatrix}.
\end{equation}
The involution induced on the space of \(\boldsymbol\theta\) is
\begin{equation}
\label{rm:eq:modular-Rckm}
\boxed{
\mathcal R_{\rm CKM}
=M_{\rm CKM}D_hM_{\rm CKM}^{-1}
=I-2c_h\ell_h.}
\end{equation}
Satisfies
\begin{equation}
\label{rm:eq:modular-Rckm-properties}
\mathcal R_{\rm CKM}^2=I,
\qquad
\det\mathcal R_{\rm CKM}=-1,
\qquad
\mathcal R_{\rm CKM}M_{\rm CKM}=M_{\rm CKM}D_h.
\end{equation}
\end{theorem}

\begin{proof}
Multiplication of
\cref{rm:eq:modular-Mckm,rm:eq:modular-Mckm-inverse} gives the identity; in
particular, \cref{rm:eq:modular-dual-coordinate-check} proves that
\(c_h\ell_h\) is the rank-one projector onto \(\mathbb Qc_h\) along
\(\operatorname{span}\{c_A,c_\Delta\}\). The exact product
\(\ell_hc_h=1\) implies
\[
(I-2c_h\ell_h)^2
=I-4c_h\ell_h+4c_h(\ell_hc_h)\ell_h=I.
\]
The operator fixes the hyperplane
\(\ker\ell_h=\operatorname{span}\{c_A,c_\Delta\}\) pointwise and changes the sign of
\(c_h\). Its eigenvalues are therefore \(1,1,-1\); hence
\(\det\mathcal R_{\rm CKM}=-1\) and
\(\Tr\mathcal R_{\rm CKM}=1\). Finally,
\[
\mathcal R_{\rm CKM}M_{\rm CKM}
=M_{\rm CKM}D_hM_{\rm CKM}^{-1}M_{\rm CKM}
=M_{\rm CKM}D_h,
\]
which proves the entanglement identity.
\end{proof}

The explicit matrix is
\begin{equation}
\label{rm:eq:modular-Rckm-matrix}
\mathcal R_{\rm CKM}=
\begin{pmatrix}
4607/3857&1760/3857&-1500/3857\\
-150/551&199/551&300/551\\
3000/3857&7040/3857&-2143/3857
\end{pmatrix}.
\end{equation}
Although \(\mathcal R_{\rm CKM}\) is not an orthogonal reflection for the
Euclidean product of the three angular coordinates, it is for the
transported rational metric
\begin{equation}
\label{rm:eq:modular-ckm-metric}
G_{\rm CKM}=(M_{\rm CKM}^{-1})^TM_{\rm CKM}^{-1},
\qquad
\mathcal R_{\rm CKM}^TG_{\rm CKM}\mathcal R_{\rm CKM}=G_{\rm CKM}.
\end{equation}
Indeed, in the basis \((c_A,c_h,c_\Delta)\), the metric is the identity
and the involution is \(D_h\). Hence the term ``reflection'' has
precise metric content and does not merely describe that
\(\mathcal R_{\rm CKM}^2=I\).
Phase \(\delta_{\rm CP}=9A\) remains fixed under this reflection. Therefore,
\(\mathcal R_{\rm CKM}\) is not physical CP conjugation, but rather the
transport of HMT sheet inversion to mixing space.

The rational inverse of \(M_{\rm CKM}\) recovers
\begin{align}
A&=\frac{1528\theta_{12}+3380\theta_{23}+801\theta_{13}}{3857},
\label{rm:eq:modular-A-from-ckm}\\
C^*&=\frac{6(75\theta_{12}+176\theta_{23}-150\theta_{13})}{3857}.
\label{rm:eq:modular-C-from-ckm}
\end{align}
The third coordinate is also recovered without loss:
\begin{equation}
\label{rm:eq:modular-Delta-from-ckm}
\Delta
=\frac{6\theta_{12}-30\theta_{23}-12\theta_{13}}{551}.
\end{equation}
These equations are coordinate changes. They do not turn CKM into a
retrospective cause of \(A,C^*\) or of \(\gammaH\).

As a terminal recognition, the previous constructor determines
\begin{equation}
\label{rm:eq:modular-ckm-values}
(\theta_{12},\theta_{23},\theta_{13},\delta_{\rm CP})
=(13.003313589509^\circ,
2.398056157332^\circ,
0.213977382244^\circ,
65.676173123554^\circ).
\end{equation}
The rational reflection is proved before using these decimals.

\begin{keyresult}
The transition \(6\to8\) is first formalized through the inclusion
\(\iota_{6,8}:\mathcal F_6\hookrightarrow\mathcal F_8\), from which
\(90,120\) and the defect \(-1/12\) are
obtained. The six-trit word then provides the bijection \(729\leftrightarrow729\), and
\(\mathcal J_{6\to8}\) transports the two collective modes to the boundary. On
that antecedent, area and the modular Hamiltonian exactly reconstruct
the two families; thermofield-double purification determines the entanglement entropy,
and the finite first law
preserves the relative-entropy defect. In CKM, sheet inversion
becomes a rational rank-one odd reflection.
\end{keyresult}

\begin{falsifier}
The initial genealogy fails if \(\iota_{6,8}\) is not injective, if the
counts are not \(90/120\), if \(\delta_{6,8}^{\rm inc}\neq-1/12\), if
\(\beta_{729}\) is not bijective or if it is falsely presented as an
additive isomorphism. Collective transport fails if
\(\mathcal J_{6\to8}D_{\rm inc}\neq
D_{\rm area}\mathcal J_{6\to8}\). Reconstruction fails if the
determinant of the boundary observables is not thirty or if
\cref{rm:eq:modular-inverse-channels} does not recover the occupations.
The TFD fails if its partial trace is not \(\rho_A\). The first law does not
extend to finite changes without the relative entropy of
\cref{rm:eq:modular-finite-first-law}. Type
\(III_{\lambda_\partial}\) is invalidated if a channel does not persist, if there is no
compatible product or if the ratio group changes. The CKM reflection
is refuted by any of
\(\mathcal R^2\neq I\), \(\det\mathcal R\neq-1\) or
\(\mathcal RM\neq MD_h\). A formula
\(\boldsymbol\theta=f(\gammaH)\) without an intertwiner violates the typed
independence of the two functional realizations of the common antecedent.
\end{falsifier}

\begin{statusbox}
Inclusion \(H\subset O\), counts \(90/120\), two normalizations of the
defect \(-1/12\), bijection of six trits \(729\leftrightarrow729\),
collective intertwiner, scale \(a_0\), finite operators, area spectrum,
TFD, oriented defect, local and finite first law, and CKM reflection:
\status{Exact internal theorem with its explicit initial structures}.
Classification
\(III_{\lambda_\partial}\): \status{Theorem with hypotheses explicit of
persistence}. Interpretation of the thermofield-double copy as a physical exterior:
\status{Realization, Physically Conditioned}. The magnitude \(\gammaH\) is taken
from the chapter on Barbero--Immirzi and minimal area, where it was derived; here it acts
as an already established structural datum.
\end{statusbox}
